\begin{prob}
    % Write the problem here.

    Suppose $T_1(n), \ldots, T_6(n)$ are functions
    describing the runtime of six algorithms. Furthermore, suppose we
    have the following bounds on each function:
    \begin{align*}
        T_1(n) &= \Theta(n^{2.5})\\
        T_2(n) &= O(n \log n)\\
        T_3(n) &= \Omega(\log n)\\
        T_4(n) &= O(n^4) \text{ and } T_4 = \Omega(n^2)\\
        T_5(n) &= \Theta(n)\\
        T_6(n) &= \Theta(n \log{n})\\
        T_7(n) &= O(n^{1.5} \log n) \text{ and } T_7 = \Omega(n \log{n})
    \end{align*}

    What are the best bounds that can be placed on the following functions?

    For this problem, you do not need to show work.

    Hint: watch the supplemental lecture at \url{https://youtu.be/tmR-bIN2qw4}.

    \textbf{Example 1}: $T_1(n) + T_2(n)$.

    \textbf{Solution}: $T_1(n) + T_2(n)$ is $\Theta(n^{2.5})$.
    
    \textbf{Example 2}: $f(n) = 2 \cdot T_4(n)$.

    \textbf{Solution}: $f(n) = 2 \cdot T_4(n)$ is $O(n^4)$ and $\Omega(n^2)$.


    % the below function looks for a file in the "solution" directory named "solution"
    % and includes it (if the solution document is being compiled)
    % \inputsolution{solution}

    \begin{subprobset}

        \begin{subprob}
            $T_1(n) + T_5(n)$

            \begin{soln}
                $\Theta(n^{2.5})$
            \end{soln}
        \end{subprob}

        \begin{subprob}
            $T_2(n) + T_6(n)$

            \begin{soln}
                $\Theta(n \log n)$
            \end{soln}
        \end{subprob}

        \begin{subprob}
            $T_2(n) + T_4(n)$

            \begin{soln}
                $O(n^4)$ and $\Omega(n^2)$
            \end{soln}
        \end{subprob}

        \begin{subprob}
            $T_7(n) + T_4(n)$

            \begin{soln}
                $O(n^4)$ and $\Omega(n^2)$
            \end{soln}
        \end{subprob}

        \begin{subprob}
            $T_3(n) + T_1(n)$

            \begin{soln}
                $\Omega(n^{2.5})$

                Note that we can't give an upper bound here, because we don't have an upper
                bound on $T_3(n)$. For example, it could be that $T_3(n) = n^{10}$ or
                $n^{100}$ or $n^{1000}$; we just do not have enough information to know.
            \end{soln}
        \end{subprob}

        \begin{subprob}
            $T_1(n) \times T_4(n)$

            \begin{soln}
                $O(n^{6.5})$ and $\Omega(n^{4.5})$

            \end{soln}
        \end{subprob}

        \begin{subprob}
            $T_6(n) + T_4(n) / T_5(n)$

            \begin{soln}
                $O(n^{3})$ and $\Omega(n \log n)$

            \end{soln}
        \end{subprob}

    \end{subprobset}

\end{prob}

